% Inactive modular snapshot; edit the root neurips_2026.tex instead.
\section{Computational cost of the spectral instruments}
\label{app:cost}
We distinguish trainable capacity from frozen storage, setup cost, and the computation required to apply the adapter. Let a square layer have width $d$, selected tail dimension $r$, and rotated sub-block dimension $q\le r$. Let $b$ count the token vectors processed by that layer in one forward call. All costs below are per adapted matrix and exclude the base model's own storage and forward computation, the task head, and hardware-dependent constants.

\paragraph{SVD initialization and frozen storage.}
A dense SVD of a $d\times d$ matrix has conventional $O(d^3)$ arithmetic cost. For an $m\times n$ matrix with $p=\min(m,n)$, the corresponding economy-SVD cost is $O(mnp)$ and its two bases contain $(m+n)p$ entries. This is a one-time setup cost per distinct weight decomposition, not a decomposition at every optimization step. The practical implementation computes the SVD in float32 before casting stored factors to the adapter dtype, so initialization workspace can exceed persistent basis storage.

A tail-only implementation needs $2dr$ basis entries to apply its update. The implemented leading-plus-tail construction retains the complete bases, totaling $2d^2$ entries, plus $O(d)$ fixed coefficients and any orthogonal-map state. Freezing these entries avoids gradients and optimizer moments for them; it does not remove their memory cost. Setting the leading coefficients to zero in a full-basis implementation does not itself discard those bases or shorten its matrix products. Diagnostics that retain an additional full reference must also be included in memory accounting.

\paragraph{Trainable entries versus effective degrees of freedom.}
Fixed-tail coefficients require $r$ trainable entries. Two dense $q\times q$ rotation parameters add $2q^2$ stored trainable entries; their orthogonal outputs have only $q(q-1)$ combined manifold degrees of freedom. Ambient input/output scalers add $2d$ entries. Thus the practical unrotated and partially rotated forms store $r+2d$ and $r+2d+2q^2$ trainable entries, respectively, before the task head. A directly trained unrestricted tail core instead has $r^2$ entries. These counts describe different parameterizations, not equal-capacity interventions. For example, $d=768,r=256$ gives 1,792 trainable entries per unrotated coordinate-scaled matrix but 1,179,648 frozen basis entries. Aggregate benchmark parameter counts in Appendix~\ref{app:downstream} are retained as reported, rather than substituted for a tensor-level memory inventory.

\paragraph{Forward computation.}
Applying a diagonal tail update through its thin factors costs $O(bdr)$; ambient scaling adds only $O(bd)$. If partial rotations are materialized into the basis first, the two basis products add $O(dq^2)$ per reconstruction, in addition to the orthogonal-map construction cost. Dense $q\times q$ maps can incur $O(q^3)$ arithmetic. A dense tail core applied between thin basis products instead costs $O(bdr+br^2)$.

The practical forward path concatenates the leading and tail factors and applies two full-basis products. Its adapter branch therefore costs $O(bd^2)$ even when $q=0$, plus rotation construction when enabled. This is additional to the $O(bd^2)$ base projection. A rank-$\ell$ LoRA branch costs $O(bd\ell)$, so a smaller trainable count need not mean a cheaper forward pass. Backpropagation also differentiates through the adapter operations, and end-to-end peak memory depends on retained activations, dtype, optimizer state, and checkpointing. The arithmetic counts alone do not establish a measured throughput ratio.

\paragraph{Regularization and diagnostics.}
Directly forming $M_E=U_L^\top E U_T$ and $M_D=V_L^\top D V_T$ costs $O(dkr)$ for a leading dimension $k=d-r$; evaluating their squared Frobenius penalties then costs $O(kr)$. These operations and their gradients are additional training work when the penalty is enabled. They do not require an SVD of the adapted weight. By contrast, exact ranked-subspace diagnostics use a new dense SVD of $W_0+\Delta$ and can cost $O(d^3)$ per measured layer/checkpoint. Diagnostic time should be reported separately from optimization time and should not be charged only to the spectral methods in a controlled comparison.

\paragraph{Merging and deployment.}
At evaluation with dropout disabled, the additive update can be materialized and merged into $W_0$. A thin diagonal-tail construction costs $O(d^2r)$ to materialize, whereas the current full-basis product costs $O(d^3)$, plus the rotation and scaler operations; dense materialization requires $O(d^2)$ output storage. After merging and removing the adapter branch, the linear layer has the base model's shape and inference-time matrix multiplication. This removes the extra branch at deployment, not the preceding SVD, training, or merge costs.
